% !TEX program = tectonic
% =============================================================================
%  AULA 8 — IoT sem Fio e LPWAN
%  Curso: Redes sem Fio  |  Profa. Anelise Munaretto
%  Conteúdo: LoRaWAN, Sigfox, NB-IoT, LTE-M, protocolos de aplicação.
%  Compilar: tectonic aula8.tex
% =============================================================================
\documentclass[aspectratio=169,11pt]{beamer}
% =============================================================================
%  PREÁMBULO COMÚN para as aulas de Redes sem Fio (Beamer + TikZ)
%  Incluido com % =============================================================================
%  PREÁMBULO COMÚN para as aulas de Redes sem Fio (Beamer + TikZ)
%  Incluido com % =============================================================================
%  PREÁMBULO COMÚN para as aulas de Redes sem Fio (Beamer + TikZ)
%  Incluido com \input{preamble.tex} en cada aula.
%  Paleta de cores e estilos são definidos aqui uma única vez.
% =============================================================================

\usepackage[T1]{fontenc}
\usepackage{amsmath,amssymb,mathtools}
\usepackage{graphicx}
\usepackage{tikz}
\usetikzlibrary{positioning,arrows.meta,calc,backgrounds,decorations.pathmorphing,fit,shadows,shapes.symbols,shapes.geometric}
\usepackage{pgfplots}
\pgfplotsset{compat=1.16}
\usepackage{tabularx,booktabs,multirow,array}
\usepackage[normalem]{ulem}
\usepackage{hyperref}

% ---------- Paleta de cores própria ----------
\definecolor{aneliseAzul}{RGB}{20,60,110}
\definecolor{aneliseVerde}{RGB}{30,140,70}
\definecolor{aneliseLaranja}{RGB}{215,120,20}
\definecolor{aneliseGris}{RGB}{90,90,90}

% ---------- Tema ----------
\usetheme{Madrid}
\usecolortheme{whale}
\setbeamercolor{structure}{fg=aneliseAzul}
\setbeamercolor{title}{fg=aneliseAzul}
\setbeamercolor{frametitle}{fg=aneliseAzul}
\setbeamercolor{block title}{fg=aneliseAzul,bg=aneliseGris!12!white}
\setbeamercolor{block body}{bg=aneliseGris!7!white}
\hypersetup{colorlinks=true,linkcolor=aneliseAzul,urlcolor=aneliseVerde}

% ---------- Comandos auxiliares ----------
\newcommand{\kurose}{Kurose \& Ross, \emph{Computer Networking: a Top-Down Approach}}
\newcommand{\forouzan}{Forouzan, \emph{Data Communications and Networking}}
\newcommand{\stallings}{Stallings, \emph{Wireless Communications \& Networks}}
\newcommand{\tanenbaum}{Tanenbaum, \emph{Computer Networks}}
\newcommand{\rfc}[1]{\href{https://www.rfc-editor.org/rfc/rfc#1.txt}{RFC~#1}}
% -----------------------------------------------------------------------------
%  Estilos TikZ comuns para desenhar redes (posicionamento relativo)
% -----------------------------------------------------------------------------
\tikzset{
  host/.style={draw,rounded corners=2mm,fill=aneliseAzul!15,inner sep=1.5mm,font=\scriptsize},
  rt/.style={draw,rounded corners=1mm,fill=aneliseLaranja!25,inner sep=1mm,font=\scriptsize},
  nube/.style={draw,cloud,aspect=2.2,fill=aneliseGris!15,inner sep=1.5mm,font=\scriptsize},
  el/.style={draw,rounded corners=1.2mm,fill=aneliseGris!18,inner sep=0.8mm,font=\scriptsize},
  enl/.style={thick},
  enlA/.style={thick,-Stealth},
} en cada aula.
%  Paleta de cores e estilos são definidos aqui uma única vez.
% =============================================================================

\usepackage[T1]{fontenc}
\usepackage{amsmath,amssymb,mathtools}
\usepackage{graphicx}
\usepackage{tikz}
\usetikzlibrary{positioning,arrows.meta,calc,backgrounds,decorations.pathmorphing,fit,shadows,shapes.symbols,shapes.geometric}
\usepackage{pgfplots}
\pgfplotsset{compat=1.16}
\usepackage{tabularx,booktabs,multirow,array}
\usepackage[normalem]{ulem}
\usepackage{hyperref}

% ---------- Paleta de cores própria ----------
\definecolor{aneliseAzul}{RGB}{20,60,110}
\definecolor{aneliseVerde}{RGB}{30,140,70}
\definecolor{aneliseLaranja}{RGB}{215,120,20}
\definecolor{aneliseGris}{RGB}{90,90,90}

% ---------- Tema ----------
\usetheme{Madrid}
\usecolortheme{whale}
\setbeamercolor{structure}{fg=aneliseAzul}
\setbeamercolor{title}{fg=aneliseAzul}
\setbeamercolor{frametitle}{fg=aneliseAzul}
\setbeamercolor{block title}{fg=aneliseAzul,bg=aneliseGris!12!white}
\setbeamercolor{block body}{bg=aneliseGris!7!white}
\hypersetup{colorlinks=true,linkcolor=aneliseAzul,urlcolor=aneliseVerde}

% ---------- Comandos auxiliares ----------
\newcommand{\kurose}{Kurose \& Ross, \emph{Computer Networking: a Top-Down Approach}}
\newcommand{\forouzan}{Forouzan, \emph{Data Communications and Networking}}
\newcommand{\stallings}{Stallings, \emph{Wireless Communications \& Networks}}
\newcommand{\tanenbaum}{Tanenbaum, \emph{Computer Networks}}
\newcommand{\rfc}[1]{\href{https://www.rfc-editor.org/rfc/rfc#1.txt}{RFC~#1}}
% -----------------------------------------------------------------------------
%  Estilos TikZ comuns para desenhar redes (posicionamento relativo)
% -----------------------------------------------------------------------------
\tikzset{
  host/.style={draw,rounded corners=2mm,fill=aneliseAzul!15,inner sep=1.5mm,font=\scriptsize},
  rt/.style={draw,rounded corners=1mm,fill=aneliseLaranja!25,inner sep=1mm,font=\scriptsize},
  nube/.style={draw,cloud,aspect=2.2,fill=aneliseGris!15,inner sep=1.5mm,font=\scriptsize},
  el/.style={draw,rounded corners=1.2mm,fill=aneliseGris!18,inner sep=0.8mm,font=\scriptsize},
  enl/.style={thick},
  enlA/.style={thick,-Stealth},
} en cada aula.
%  Paleta de cores e estilos são definidos aqui uma única vez.
% =============================================================================

\usepackage[T1]{fontenc}
\usepackage{amsmath,amssymb,mathtools}
\usepackage{graphicx}
\usepackage{tikz}
\usetikzlibrary{positioning,arrows.meta,calc,backgrounds,decorations.pathmorphing,fit,shadows,shapes.symbols,shapes.geometric}
\usepackage{pgfplots}
\pgfplotsset{compat=1.16}
\usepackage{tabularx,booktabs,multirow,array}
\usepackage[normalem]{ulem}
\usepackage{hyperref}

% ---------- Paleta de cores própria ----------
\definecolor{aneliseAzul}{RGB}{20,60,110}
\definecolor{aneliseVerde}{RGB}{30,140,70}
\definecolor{aneliseLaranja}{RGB}{215,120,20}
\definecolor{aneliseGris}{RGB}{90,90,90}

% ---------- Tema ----------
\usetheme{Madrid}
\usecolortheme{whale}
\setbeamercolor{structure}{fg=aneliseAzul}
\setbeamercolor{title}{fg=aneliseAzul}
\setbeamercolor{frametitle}{fg=aneliseAzul}
\setbeamercolor{block title}{fg=aneliseAzul,bg=aneliseGris!12!white}
\setbeamercolor{block body}{bg=aneliseGris!7!white}
\hypersetup{colorlinks=true,linkcolor=aneliseAzul,urlcolor=aneliseVerde}

% ---------- Comandos auxiliares ----------
\newcommand{\kurose}{Kurose \& Ross, \emph{Computer Networking: a Top-Down Approach}}
\newcommand{\forouzan}{Forouzan, \emph{Data Communications and Networking}}
\newcommand{\stallings}{Stallings, \emph{Wireless Communications \& Networks}}
\newcommand{\tanenbaum}{Tanenbaum, \emph{Computer Networks}}
\newcommand{\rfc}[1]{\href{https://www.rfc-editor.org/rfc/rfc#1.txt}{RFC~#1}}
% -----------------------------------------------------------------------------
%  Estilos TikZ comuns para desenhar redes (posicionamento relativo)
% -----------------------------------------------------------------------------
\tikzset{
  host/.style={draw,rounded corners=2mm,fill=aneliseAzul!15,inner sep=1.5mm,font=\scriptsize},
  rt/.style={draw,rounded corners=1mm,fill=aneliseLaranja!25,inner sep=1mm,font=\scriptsize},
  nube/.style={draw,cloud,aspect=2.2,fill=aneliseGris!15,inner sep=1.5mm,font=\scriptsize},
  el/.style={draw,rounded corners=1.2mm,fill=aneliseGris!18,inner sep=0.8mm,font=\scriptsize},
  enl/.style={thick},
  enlA/.style={thick,-Stealth},
}

\title[Redes sem Fio]{Redes sem Fio\\ IoT sem Fio e LPWAN}
\subtitle{Aula 8}
\author[Anelise Munaretto]{Profa. Anelise Munaretto}
\institute[Curso]{Redes sem Fio --- Ciência da Computação}
\date{2026}

\begin{document}

\begin{frame}[plain]
  \titlepage
\end{frame}

\section{Introdução}

\begin{frame}{IoT --- Internet das Coisas}
  \begin{block}{Definição}
    Rede de objetos físicos (\emph{coisas}) equipados com sensores, atuadores e conectividade, capazes de coletar e trocar dados.
  \end{block}
  \begin{itemize}
    \item Projeção: 30--50 bilhões de dispositivos IoT até 2030
    \item Desafios: bateria, alcance, custo, densidade, segurança
  \end{itemize}
  \begin{columns}
    \begin{column}{0.5\textwidth}
      \textbf{Tecnologias LPWAN}
      \begin{itemize}
        \item LoRaWAN
        \item Sigfox
        \item NB-IoT / LTE-M
        \item Wi-SUN, mioty
      \end{itemize}
    \end{column}
    \begin{column}{0.5\textwidth}
      \textbf{Atributos típicos}
      \begin{itemize}
        \item alcance: 2--15~km (rural até 50~km)
        \item taxa: 0,1--250~kbps
        \item bateria: 5--10 anos
        \item custo de módulo: \$2--\$10
      \end{itemize}
    \end{column}
  \end{columns}
\end{frame}

\section{LoRaWAN}

\begin{frame}{LoRaWAN --- arquitetura e operação}
  \begin{columns}
    \begin{column}{0.55\textwidth}
      \textbf{Camada física: LoRa (CSS)}
      \begin{itemize}
        \item Chirp Spread Spectrum (sub-GHz: 868/915~MHz)
        \item 6 \emph{spreading factors} (SF7--SF12)
        \item sensibilidade até $-137$~dBm
      \end{itemize}
      \textbf{Camada MAC: LoRaWAN}
      \begin{itemize}
        \item topologia \emph{star-of-stars}
        \item \textbf{Classes}: A (ALOHA), B (beacon), C (escuta contínua)
        \item \emph{Adaptive Data Rate} (ADR)
        \item autenticação AES-128 (Join-Request/Accept)
      \end{itemize}
    \end{column}
    \begin{column}{0.45\textwidth}
      \begin{center}
        \begin{tikzpicture}[font=\scriptsize,
            nd/.style={draw,circle,minimum size=5mm,fill=aneliseAzul!15,inner sep=0.4mm},
            gw/.style={draw,rounded corners=1.5mm,fill=aneliseLaranja!25,inner sep=1mm,text width=1.8cm,align=center},
          ]
          \node[nd](s1){S};
          \node[nd,right=0.8cm of s1](s2){S};
          \node[nd,right=0.8cm of s2](s3){S};
          \node[gw,above=1.4cm of s2](gw1){Gateway};
          \node[gw,right=1.4cm of gw1](gw2){Gateway};
          \node[draw,fill=aneliseVerde!20,rounded corners=1.5mm,above=1.4cm of gw2,text width=2.2cm,align=center](ns){Servidor\\de rede};
          \draw[thick,aneliseVerde] (s1) -- (gw1);
          \draw[thick,aneliseVerde] (s2) -- (gw1);
          \draw[thick,aneliseVerde] (s3) -- (gw2);
          \draw[thick,aneliseLaranja] (gw1) -- (gw2);
          \draw[thick,aneliseLaranja] (gw2) -- (ns);
          \node[below=1mm of s1,font=\tiny]{\emph{end-device}};
        \end{tikzpicture}
      \end{center}
    \end{column}
  \end{columns}
\end{frame}

\begin{frame}{LoRaWAN: classes e ADR}
  \begin{itemize}
    \item \textbf{Classe A}: uplink seguido de duas janelas de downlink (mínimo consumo)
    \item \textbf{Classe B}: janelas de recepção periódicas (beacon)
    \item \textbf{Classe C}: receptor sempre ligado (downlink contínuo)
  \end{itemize}
  \begin{block}{ADR --- Adaptive Data Rate}
    \begin{itemize}
      \item A rede otimiza SF e potência de transmissão
      \item SF12 $\rightarrow$ maior alcance, menor taxa (290~bps)
      \item SF7 $\rightarrow$ menor alcance, maior taxa (5,4~kbps)
      \item Economia de bateria (até 10 anos com 2 pilhas AA)
    \end{itemize}
  \end{block}
\end{frame}

\section{Sigfox}

\begin{frame}{Sigfox --- UNB (Ultra-Narrow Band)}
  \begin{columns}
    \begin{column}{0.55\textwidth}
      \begin{itemize}
        \item Banda ISM sub-GHz (868/915~MHz)
        \item Largura de canal: 100~Hz (ultra-estreita)
        \item Uplink: 140 mensagens/dia, 100~bps
        \item Downlink: 4 mensagens/dia (opcional)
        \item Alcance: 10--50~km (rural)
        \item Modelo de operadora (SIGFOX)
      \end{itemize}
    \end{column}
    \begin{column}{0.45\textwidth}
      \begin{block}{Comparativo LoRaWAN vs.\ Sigfox}
        \begin{itemize}
          \item LoRaWAN: bidirecional, flexível, redes privadas
          \item Sigfox: ultra-simples, unidirecional, operadora
        \end{itemize}
      \end{block}
    \end{column}
  \end{columns}
\end{frame}

\section{NB-IoT e LTE-M}

\begin{frame}{NB-IoT e LTE-M (3GPP)}
  \begin{columns}
    \begin{column}{0.5\textwidth}
      \textbf{NB-IoT (Narrowband IoT)}
      \begin{itemize}
        \item 200~kHz, dentro da banda LTE
        \item taxa: até 250~kbps (DL), 20~kbps (UL)
        \item penetração indoor superior
        \item sem mobilidade (fixo/estacionário)
      \end{itemize}
    \end{column}
    \begin{column}{0.5\textwidth}
      \textbf{LTE-M (eMTC)}
      \begin{itemize}
        \item 1,4~MHz, dentro da banda LTE
        \item taxa: até 1~Mbps
        \item suporta mobilidade e \emph{voice}
        \item ideal para rastreadores, wearables
      \end{itemize}
    \end{column}
  \end{columns}
  \vspace{0.2cm}
  \begin{block}{Vantagens sobre LPWAN não-celular}
    \begin{itemize}
      \item QoS garantida (recurso dedicado)
      \item segurança nativa (SIM, AKA)
      \item integração com 5G core (mMTC)
    \end{itemize}
  \end{block}
\end{frame}

\section{Protocolos de aplicação}

\begin{frame}{Protocolos de aplicação para IoT}
  \begin{columns}
    \begin{column}{0.5\textwidth}
      \textbf{MQTT}
      \begin{itemize}
        \item \emph{Subscribe/Publish} via broker
        \item leve, suporta QoS 0/1/2
        \item padrão de facto para IoT
      \end{itemize}
      \textbf{CoAP}
      \begin{itemize}
        \item UDP + REST (como HTTP)
        \item multicast nativo
        \item muito usado em Constrained RESTful Environments
      \end{itemize}
    \end{column}
    \begin{column}{0.5\textwidth}
      \textbf{LwM2M (Lightweight M2M)}
      \begin{itemize}
        \item gerenciamento de dispositivo
        \item sobre CoAP ou MQTT
        \item padrão OMA SpecWorks
      \end{itemize}
      \textbf{HTTP/2 e gRPC}
      \begin{itemize}
        \item para gateways/pontos de agregação
        \item streaming bidirecional
      \end{itemize}
    \end{column}
  \end{columns}
\end{frame}

\section{Bibliografia}

\begin{frame}{Bibliografia}
  \begin{itemize}
    \item LoRa Alliance, \emph{LoRaWAN 1.0.4 Specification}, 2020
    \item 3GPP TS 36.300, \emph{Evolved Universal Terrestrial Radio Access (E-UTRA) and Evolved Packet Core (EPC)}
    \item 3GPP TR 45.820, \emph{Cellular system support for ultra-low complexity and low throughput Internet of Things (CIoT)}
    \item J. de Carvalho Silva et al., ``LoRaWAN --- A low power WAN protocol for IoT: A review and opportunities'', \emph{Computer Networks}, 2017
    \item MQTT v5.0 OASIS Standard, 2019
  \end{itemize}
\end{frame}

\end{document}